% page 188 of the facsimile (image p-187 of the scan)
% block 160.0 x 242.0 mm ; left margin 8.0 mm
\pgc{-12.700mm}{0.000mm}{
\l{\cel{3}180}
\l{}
\l{}
\l{\sou{frikasar}. (\sou{trans.}) Koquar en sauco (legumi, o karno pecig-}
\l{\cel{3}ita). - DEFIRS.}
\l{}
\l{\sou{frikaseo.}\cel{1}Raguto de karno pecigita, e koquita en sauco.}
\l{\cel{3}DEFIRS.}
\l{}
\l{\sou{fripo}. Vesti, linjo, mobli olda, de qui onu su desembarasas}
\l{\cel{3}per vendar li a brokantisti. - EF.}
\l{}
\l{\sou{fripono}. I. Persono qua habile furtas mikra kozi. - II. Per-}
\l{\cel{3}sono qua agas mala stroki malicoze. - F.}
\l{}
\l{\sou{friso}. I. En la monumenti di stilo Greka : parto di la en-}
\l{\cel{3}tablamento inter la arkitravo e la kornico. - II. Bendo}
\l{\cel{3}piktita o skultita cirkum la parto supra di chambro, di}
\l{\cel{3}vazo, super la kadro di pordo, di kameno, e c.-DEFIRS.}
\l{}
\l{\sou{friskar}. (\sou{netrans.}) Agitar (su) per movi vivaca e kurta. - E.}
\l{}
\l{\sou{fristo}. Tempo-quanto grantenda o grantita por efektigar}
\l{\cel{3}ulo. - D.}
\l{}
\l{\sou{fritar}.\cel{2}(\sou{trans}.)Koquar per padelo, en graso, oleo od en}
\l{\cel{3}butro bolianta. - EFIS.}
\l{}
\l{\sou{fritilario}. (\sou{bot}.)\cel{2}Planto liliacea, di qua la floro havas}
\l{\cel{3}la formo di korneto subversita, e di qua la bulbo pulpoza}
\l{\cel{3}kontenas elemento akra, drastika. - L.\cel{1}\sou{fritilius}.}
\l{}
\l{\sou{frivola}. I. Tante vana ke lu ne meritas ke onu fervorez por}
\l{\cel{3}lu. - II. Qua fervoras por la kozi vana. - DEFIRS.}
\l{}
\l{\sou{frizar}. (\sou{trans.}) I. Volvar (hari, pili, e c.) cirkum li, preske}
\l{\cel{3}lokle. - II. (\sou{netrans}.)(\sou{aludante hari, pili, e}\cel{1}c.) Ipse}
\l{\cel{3}volvar cirkum su, preske lokle. - DEFS.}
\l{}
\l{\sou{frogo}. Sorto di kazako kun longa maniki, di qua la uzo-kustumo}
\l{\cel{3}venis de Brandeburgia, lor la rejo Ludovikus XIV. di Fran-}
\l{\cel{3}cia, ordinare garnisita ye butoni olivatra, interligita}
\l{\cel{3}per galoni. - F.}
\l{}
\l{\sou{froko.}\cel{2}Vesto di monako katolika. - EF.}
\l{}
\l{\sou{frolar}.\cel{2}(\sou{trans.}) Tushar lejere la bordo, la extremajo di}
\l{\cel{3}ulo. - F.}
\l{}
\l{\sou{fromajo}. Substanco nutrala quan onu obtenas de preparar di-}
\l{\cel{3}verse butro, la parto kazeatra di lakto, o ta du materii}
\l{\cel{3}kombinita. - FI.}
\l{}
\l{frondo.\cel{2}Lans-armo, posho, quan onu igas jirar per du kordi}
\l{\cel{3}qui lu suspendesas, e qua kontenas stono o globeto quan}
\l{\cel{3}onu lasas eskapar bruske. - FI.}
}
